%% Generated by Sphinx.
\def\sphinxdocclass{report}
\documentclass[a4paper,10pt,english]{sphinxmanual}
\ifdefined\pdfpxdimen
   \let\sphinxpxdimen\pdfpxdimen\else\newdimen\sphinxpxdimen
\fi \sphinxpxdimen=.75bp\relax
\ifdefined\pdfimageresolution
    \pdfimageresolution= \numexpr \dimexpr1in\relax/\sphinxpxdimen\relax
\fi
\newdimen\sphinxremdimen\sphinxremdimen = 10pt
%% let collapsible pdf bookmarks panel have high depth per default
\PassOptionsToPackage{bookmarksdepth=5}{hyperref}

\PassOptionsToPackage{booktabs}{sphinx}
\PassOptionsToPackage{colorrows}{sphinx}

\PassOptionsToPackage{warn}{textcomp}
\usepackage[utf8]{inputenc}
\ifdefined\DeclareUnicodeCharacter
% support both utf8 and utf8x syntaxes
  \ifdefined\DeclareUnicodeCharacterAsOptional
    \def\sphinxDUC#1{\DeclareUnicodeCharacter{"#1}}
  \else
    \let\sphinxDUC\DeclareUnicodeCharacter
  \fi
  \sphinxDUC{00A0}{\nobreakspace}
  \sphinxDUC{2500}{\sphinxunichar{2500}}
  \sphinxDUC{2502}{\sphinxunichar{2502}}
  \sphinxDUC{2514}{\sphinxunichar{2514}}
  \sphinxDUC{251C}{\sphinxunichar{251C}}
  \sphinxDUC{2572}{\textbackslash}
\fi
\usepackage{cmap}
\usepackage[T1]{fontenc}
\usepackage{amsmath,amssymb,amstext}
\usepackage{babel}



\usepackage{tgtermes}
\usepackage{tgheros}
\renewcommand{\ttdefault}{txtt}



\usepackage[Bjarne]{fncychap}
\usepackage[,numfigreset=1,mathnumfig,mathnumsep={.}]{sphinx}

\fvset{fontsize=auto}
\usepackage{geometry}

\usepackage{nbsphinx}

% Include hyperref last.
\usepackage{hyperref}
% Fix anchor placement for figures with captions.
\usepackage{hypcap}% it must be loaded after hyperref.
% Set up styles of URL: it should be placed after hyperref.
\urlstyle{same}

\addto\captionsenglish{\renewcommand{\contentsname}{Contents}}

\usepackage{sphinxmessages}
\setcounter{tocdepth}{3}
\setcounter{secnumdepth}{3}

\usepackage{xcolor}
\usepackage[
    type={CC},
    modifier={by},
    version={4.0},
]{doclicense}

\newcommand{\docID}{ESA-RL2OCEAN-CALVAL-ATBD}
\newcommand{\docversion}{0.2}
\newcommand{\doctitle}{RL2Ocean Calibration and Validation ATBD}

\fancypagestyle{normal}{
  \fancyhf{}
  \fancyhead[L]{\color{gray}\docID}
  \fancyhead[C]{\color{gray}\doctitle}
  \fancyhead[R]{\color{gray}\docversion}
  \fancyfoot[C]{\doclicenseImage[imagewidth=2cm]}
  \fancyfoot[R]{\small\thepage}
  \renewcommand{\headrulewidth}{0.0pt}
  \renewcommand{\footrulewidth}{0.0pt}
}
\fancypagestyle{plain}{
  \fancyhf{}
  \fancyhead[L]{\color{gray}\docID}
  \fancyhead[C]{\color{gray}\doctitle}
  \fancyhead[R]{\color{gray}\docversion}
  \fancyfoot[C]{\doclicenseImage[imagewidth=2cm]}
  \fancyfoot[R]{\small\thepage}
  \renewcommand{\headrulewidth}{0.0pt}
  \renewcommand{\footrulewidth}{0.0pt}
}
\fancypagestyle{empty}{
  \fancyhf{}
  \fancyfoot[C]{\doclicenseImage[imagewidth=2cm]}
  \renewcommand{\headrulewidth}{0pt}
  \renewcommand{\footrulewidth}{0pt}
}


\title{RL2Ocean Calibration and Validation ATBD}
\date{2026-09-29}
\release{0.2}
\author{RL2Ocean consortium}
\newcommand{\sphinxlogo}{\vbox{}}
\renewcommand{\releasename}{Release}
\makeindex
\begin{document}

\ifdefined\shorthandoff
  \ifnum\catcode`\=\string=\active\shorthandoff{=}\fi
  \ifnum\catcode`\"=\active\shorthandoff{"}\fi
\fi

\pagestyle{empty}
\sphinxmaketitle
\pagestyle{plain}
\sphinxtableofcontents
\pagestyle{normal}
\phantomsection\label{\detokenize{index::doc}}


\noindent{\hspace*{\fill}\sphinxincludegraphics[width=0.850\linewidth]{{partners_logo}.png}\hspace*{\fill}}


\begin{savenotes}\sphinxattablestart
\sphinxthistablewithglobalstyle
\centering
\sphinxcapstartof{table}
\sphinxthecaptionisattop
\sphinxcaption{Document information}\label{\detokenize{index:id1}}
\sphinxaftertopcaption
\begin{tabular}[t]{\X{30}{100}\X{70}{100}}
\sphinxtoprule
\sphinxtableatstartofbodyhook\begin{varwidth}[t]{\sphinxcolwidth{1}{2}}
\sphinxAtStartPar
Document ID
\sphinxbeforeendvarwidth
\end{varwidth}%
&\begin{varwidth}[t]{\sphinxcolwidth{1}{2}}
\sphinxAtStartPar
ESA\sphinxhyphen{}RL2OCEAN\sphinxhyphen{}CALVAL\sphinxhyphen{}ATBD
\sphinxbeforeendvarwidth
\end{varwidth}%
\\
\sphinxhline\begin{varwidth}[t]{\sphinxcolwidth{1}{2}}
\sphinxAtStartPar
Issue
\sphinxbeforeendvarwidth
\end{varwidth}%
&\begin{varwidth}[t]{\sphinxcolwidth{1}{2}}
\sphinxAtStartPar
0.2
\sphinxbeforeendvarwidth
\end{varwidth}%
\\
\sphinxhline\begin{varwidth}[t]{\sphinxcolwidth{1}{2}}
\sphinxAtStartPar
Date
\sphinxbeforeendvarwidth
\end{varwidth}%
&\begin{varwidth}[t]{\sphinxcolwidth{1}{2}}
\sphinxAtStartPar
2026\sphinxhyphen{}09\sphinxhyphen{}29
\sphinxbeforeendvarwidth
\end{varwidth}%
\\
\sphinxhline\begin{varwidth}[t]{\sphinxcolwidth{1}{2}}
\sphinxAtStartPar
Status
\sphinxbeforeendvarwidth
\end{varwidth}%
&\begin{varwidth}[t]{\sphinxcolwidth{1}{2}}
\sphinxAtStartPar
Draft (approval status not stated in source)
\sphinxbeforeendvarwidth
\end{varwidth}%
\\
\sphinxhline\begin{varwidth}[t]{\sphinxcolwidth{1}{2}}
\sphinxAtStartPar
Processor
\sphinxbeforeendvarwidth
\end{varwidth}%
&\begin{varwidth}[t]{\sphinxcolwidth{1}{2}}
\sphinxAtStartPar
ROSE\sphinxhyphen{}L L2 Ocean Processor (RL2Ocean)
\sphinxbeforeendvarwidth
\end{varwidth}%
\\
\sphinxhline\begin{varwidth}[t]{\sphinxcolwidth{1}{2}}
\sphinxAtStartPar
Copyright
\sphinxbeforeendvarwidth
\end{varwidth}%
&\begin{varwidth}[t]{\sphinxcolwidth{1}{2}}
\sphinxAtStartPar
RL2Ocean consortium
\sphinxbeforeendvarwidth
\end{varwidth}%
\\
\sphinxhline\begin{varwidth}[t]{\sphinxcolwidth{1}{2}}
\sphinxAtStartPar
Distribution and license
\sphinxbeforeendvarwidth
\end{varwidth}%
&\begin{varwidth}[t]{\sphinxcolwidth{1}{2}}
\sphinxAtStartPar
\sphinxhref{https://creativecommons.org/licenses/by/4.0}{CC BY 4.0}%
\begin{footnote}[1]\sphinxAtStartFootnote
\sphinxnolinkurl{https://creativecommons.org/licenses/by/4.0}
%
\end{footnote}
\sphinxbeforeendvarwidth
\end{varwidth}%
\\
\sphinxbottomrule
\end{tabular}
\sphinxtableafterendhook\par
\sphinxattableend\end{savenotes}

\noindent{\hspace*{\fill}\sphinxincludegraphics[width=0.300\linewidth]{{ESA_logo}.png}\hspace*{\fill}}

\noindent{\hspace*{\fill}\sphinxincludegraphics[width=0.200\linewidth]{{rl2ocean_logo}.png}\hspace*{\fill}}


\chapter{Authors and approval}
\label{\detokenize{index:authors-and-approval}}
\sphinxAtStartPar
The LaTeX cover lists author placeholders rather than individual names.


\begin{savenotes}\sphinxattablestart
\sphinxthistablewithglobalstyle
\centering
\begin{tabular}[t]{\X{40}{100}\X{60}{100}}
\sphinxtoprule
\begin{varwidth}[t]{\sphinxcolwidth{1}{2}}
\sphinxstyletheadfamily \sphinxAtStartPar
Author
\sphinxbeforeendvarwidth
\end{varwidth}%
&\begin{varwidth}[t]{\sphinxcolwidth{1}{2}}
\sphinxstyletheadfamily \sphinxAtStartPar
Organisation
\sphinxbeforeendvarwidth
\end{varwidth}%
\\
\sphinxmidrule
\sphinxtableatstartofbodyhook\begin{varwidth}[t]{\sphinxcolwidth{1}{2}}
\sphinxAtStartPar
Author 1 (name not provided)
\sphinxbeforeendvarwidth
\end{varwidth}%
&\begin{varwidth}[t]{\sphinxcolwidth{1}{2}}
\sphinxAtStartPar
Norwegian Meteorological Institute (MET Norway)
\sphinxbeforeendvarwidth
\end{varwidth}%
\\
\sphinxhline\begin{varwidth}[t]{\sphinxcolwidth{1}{2}}
\sphinxAtStartPar
Author 2 (name not provided)
\sphinxbeforeendvarwidth
\end{varwidth}%
&\begin{varwidth}[t]{\sphinxcolwidth{1}{2}}
\sphinxAtStartPar
Institute of Marine Sciences (ISMAR\sphinxhyphen{}CNR)
\sphinxbeforeendvarwidth
\end{varwidth}%
\\
\sphinxhline\begin{varwidth}[t]{\sphinxcolwidth{1}{2}}
\sphinxAtStartPar
Author 3 (name not provided)
\sphinxbeforeendvarwidth
\end{varwidth}%
&\begin{varwidth}[t]{\sphinxcolwidth{1}{2}}
\sphinxAtStartPar
Barcelona Expert Center (BEC ICM\sphinxhyphen{}CSIC)
\sphinxbeforeendvarwidth
\end{varwidth}%
\\
\sphinxhline\begin{varwidth}[t]{\sphinxcolwidth{1}{2}}
\sphinxAtStartPar
Author 4 (name not provided)
\sphinxbeforeendvarwidth
\end{varwidth}%
&\begin{varwidth}[t]{\sphinxcolwidth{1}{2}}
\sphinxAtStartPar
Delft University of Technology (TU Delft)
\sphinxbeforeendvarwidth
\end{varwidth}%
\\
\sphinxhline\begin{varwidth}[t]{\sphinxcolwidth{1}{2}}
\sphinxAtStartPar
Author 5 (name not provided)
\sphinxbeforeendvarwidth
\end{varwidth}%
&\begin{varwidth}[t]{\sphinxcolwidth{1}{2}}
\sphinxAtStartPar
German Aerospace Center (DLR)
\sphinxbeforeendvarwidth
\end{varwidth}%
\\
\sphinxbottomrule
\end{tabular}
\sphinxtableafterendhook\par
\sphinxattableend\end{savenotes}

\sphinxAtStartPar
Approval and authorisation records are not specified in the source.


\chapter{Change records}
\label{\detokenize{index:change-records}}

\begin{savenotes}\sphinxattablestart
\sphinxthistablewithglobalstyle
\centering
\begin{tabular}[t]{\X{15}{100}\X{20}{100}\X{45}{100}\X{20}{100}}
\sphinxtoprule
\begin{varwidth}[t]{\sphinxcolwidth{1}{4}}
\sphinxstyletheadfamily \sphinxAtStartPar
Issue
\sphinxbeforeendvarwidth
\end{varwidth}%
&\begin{varwidth}[t]{\sphinxcolwidth{1}{4}}
\sphinxstyletheadfamily \sphinxAtStartPar
Date
\sphinxbeforeendvarwidth
\end{varwidth}%
&\begin{varwidth}[t]{\sphinxcolwidth{1}{4}}
\sphinxstyletheadfamily \sphinxAtStartPar
Description
\sphinxbeforeendvarwidth
\end{varwidth}%
&\begin{varwidth}[t]{\sphinxcolwidth{1}{4}}
\sphinxstyletheadfamily \sphinxAtStartPar
Author
\sphinxbeforeendvarwidth
\end{varwidth}%
\\
\sphinxmidrule
\sphinxtableatstartofbodyhook\begin{varwidth}[t]{\sphinxcolwidth{1}{4}}
\sphinxAtStartPar
0.2
\sphinxbeforeendvarwidth
\end{varwidth}%
&\begin{varwidth}[t]{\sphinxcolwidth{1}{4}}
\sphinxAtStartPar
Not stated
\sphinxbeforeendvarwidth
\end{varwidth}%
&\begin{varwidth}[t]{\sphinxcolwidth{1}{4}}
\sphinxAtStartPar
Source issue
\sphinxbeforeendvarwidth
\end{varwidth}%
&\begin{varwidth}[t]{\sphinxcolwidth{1}{4}}
\sphinxAtStartPar
RL2Ocean consortium
\sphinxbeforeendvarwidth
\end{varwidth}%
\\
\sphinxhline\begin{varwidth}[t]{\sphinxcolwidth{1}{4}}
\sphinxAtStartPar
0.2
\sphinxbeforeendvarwidth
\end{varwidth}%
&\begin{varwidth}[t]{\sphinxcolwidth{1}{4}}
\sphinxAtStartPar
2026\sphinxhyphen{}09\sphinxhyphen{}29
\sphinxbeforeendvarwidth
\end{varwidth}%
&\begin{varwidth}[t]{\sphinxcolwidth{1}{4}}
\sphinxAtStartPar
Sphinx migration
\sphinxbeforeendvarwidth
\end{varwidth}%
&\begin{varwidth}[t]{\sphinxcolwidth{1}{4}}
\sphinxAtStartPar
RL2Ocean consortium
\sphinxbeforeendvarwidth
\end{varwidth}%
\\
\sphinxbottomrule
\end{tabular}
\sphinxtableafterendhook\par
\sphinxattableend\end{savenotes}

\begin{sphinxadmonition}{note}{Note:}
\sphinxAtStartPar
The source abstract contains only Lorem ipsum placeholder text and is not
reproduced. The source’s substantive content and section outline are
preserved; headings without body text are identified in the relevant
chapters rather than filled with invented material.
\end{sphinxadmonition}

\sphinxstepscope


\section{Introduction}
\label{\detokenize{01_introduction:introduction}}\label{\detokenize{01_introduction::doc}}

\subsection{Purpose and scope}
\label{\detokenize{01_introduction:purpose-and-scope}}
\sphinxAtStartPar
The Copernicus Programme, established by the European Union, is a framework for Earth Observation and Monitoring, designed to support environmental protection, civil protection, and civil security efforts. At the heart of this programme is the Copernicus Space Component (CSC), which comprises the Sentinel missions. These missions are integral to providing essential data across various thematic domains, thereby facilitating informed decision\sphinxhyphen{}making and policy development.

\sphinxAtStartPar
One of the new advancements within this framework is the Copernicus Radar Observing System for Europe \textendash{} L\sphinxhyphen{}Band (ROSE\sphinxhyphen{}L). This initiative is aimed at enhancing existing capabilities by introducing L\sphinxhyphen{}band Synthetic Aperture Radar (SAR) data. The L\sphinxhyphen{}band SAR is set to complement the C\sphinxhyphen{}band radar data provided by Sentinel\sphinxhyphen{}1, thus addressing specific observational needs within environmental and maritime surveillance.


\subsection{Cal/Val context}
\label{\detokenize{01_introduction:cal-val-context}}
\sphinxAtStartPar
Quality Assessment (QA) is required for remote sensing and Earth Observation (EO) instruments in order to produce reliable measurements and ensure accuracy. Quality of products are commonly assessed through a Calibration and Validation process (Cal/Val). Consequently, the development of the prototype and operational ROSE\sphinxhyphen{}L Ocean products processors (L2O\sphinxhyphen{}GPP and L2O\sphinxhyphen{}IPF respectively) must be followed by the validation of products. A validation strategy is meant to be planned prior the calibration and validation activities. The required strategy shall cover different stages of the processor and different phases of ROSE\sphinxhyphen{}L mission, from pre\sphinxhyphen{}launch (phase D) to exploitation (phase E2). The strategy shall include methodologies, datasets, performance and uncertainty metrics.

\sphinxAtStartPar
The ROSE\sphinxhyphen{}L instrument will operate in scanning synthetic aperture radar (ScanSAR) mode. The system will support Dual\sphinxhyphen{}Pol, Quad\sphinxhyphen{}Pol and Wave modes. The Cal/Val strategy is also meant to cover all modes and polarizations. ROSE\sphinxhyphen{}L will not be the first SAR system functioning within L\sphinxhyphen{}band. ALOS, ALOS\sphinxhyphen{}2, ALOS\sphinxhyphen{}4, SAOCOM and NiSAR are, or have been, operating in L\sphinxhyphen{}band too. However, there are differences between those systems and the ROSE\sphinxhyphen{}L system, e.g. incidence angle, range bandwith, swath width, among others. Furthermore, ROSE\sphinxhyphen{}L incorporates the use of a new technique, the Scan\sphinxhyphen{}On\sphinxhyphen{}Receive (SCORE) technique with Multiple Azimuth Phase Centers (MAPS). The different instrument design will impose differences in the retrievals at low level which are meant to influence the inversion into L2 products.

\sphinxAtStartPar
Regarding the development of a retrieval algorithm but also Cal/Val activities, instrument differences between ROSE\sphinxhyphen{}L and previous SAR L\sphinxhyphen{}band instruments impose a challenge; the need of producing and validating ocean products from ROSE\sphinxhyphen{}L scenes during the pre\sphinxhyphen{}launch phase. The lack of real ROSE\sphinxhyphen{}L data will be mitigated by the use of similar acquisitions by the above\sphinxhyphen{}mentioned systems; although this introduces a non\sphinxhyphen{}negligible degree of uncertainty. For Cal/Val purposes, simulated ROSE\sphinxhyphen{}L scenes will be produced. The achievable accuracy of simulated ROSE\sphinxhyphen{}L will also be analysed during the In\sphinxhyphen{}Orbit\sphinxhyphen{}Commissioning (IOC) phase, E1. Once real ROSE\sphinxhyphen{}L data is available, the outputs from the L2O\sphinxhyphen{}IPF processor will need to be validated again, with the possibility to encounter a need for a readjustment of the processor or introduction of supplementary calibration steps.

\sphinxAtStartPar
A tool for assessing the quality of the ocean products is vital. The Cal/Val tool, L2O\sphinxhyphen{}QC, is meant to be developed in parallel to the development of the ocean processors, L2O\sphinxhyphen{}GPP and L2O\sphinxhyphen{}IPF. The L2O\sphinxhyphen{}QC software is required to be able to ingest datasets of different nature; ROSE\sphinxhyphen{}L acquisitions, reference datasets (either in\sphinxhyphen{}situ or remote), and auxiliary data. The goal is to collocate the different datasets in order to undertake the quality assessment, where the accuracy and reliability of the ROSE\sphinxhyphen{}L L2 ocean products will be determined and ensured. The underlying objective is to provide users with a set of performance and uncertainty metrics necessary for a deep understanding of the products and posterior correct usage and application.


\subsection{L2 product scope}
\label{\detokenize{01_introduction:l2-product-scope}}
\sphinxAtStartPar
The L2 products to validate are already listed in the objectives of L2O\sphinxhyphen{}GPP and L2O\sphinxhyphen{}IPF: Ocean Surface Vector Stress (OSVS), Sea State, and line\sphinxhyphen{}of\sphinxhyphen{}sight (LoS) surface currents. Within each of the three thematic products, the following L2 variables are meant to be produced and therefore validated:
\begin{itemize}
\item {} 
\sphinxAtStartPar
OSVS:
\begin{itemize}
\item {} 
\sphinxAtStartPar
Wind vector

\item {} 
\sphinxAtStartPar
Wind stress

\item {} 
\sphinxAtStartPar
Hub height wind speed

\end{itemize}

\item {} 
\sphinxAtStartPar
Sea state:
\begin{itemize}
\item {} 
\sphinxAtStartPar
Total significant wave height

\item {} 
\sphinxAtStartPar
Mean period wind\sphinxhyphen{}sea

\item {} 
\sphinxAtStartPar
First moment wave period

\item {} 
\sphinxAtStartPar
Second moment wave period

\item {} 
\sphinxAtStartPar
Wave height swell dominant system

\item {} 
\sphinxAtStartPar
Wave height swell secondary system

\item {} 
\sphinxAtStartPar
Significant wave height wind\sphinxhyphen{}sea

\item {} 
\sphinxAtStartPar
Mean period wind\sphinxhyphen{}sea

\item {} 
\sphinxAtStartPar
Wave direction

\item {} 
\sphinxAtStartPar
Wave directional energy spectrum

\end{itemize}

\item {} 
\sphinxAtStartPar
LoS currents:
\begin{itemize}
\item {} 
\sphinxAtStartPar
LoS wave Doppler

\item {} 
\sphinxAtStartPar
LoS surface current component

\end{itemize}

\end{itemize}


\subsection{Cal/Val objectives and performance}
\label{\detokenize{01_introduction:cal-val-objectives-and-performance}}
\sphinxAtStartPar
The key objectives of the Calibration and Validation (Cal/Val) phase include ensuring that all L2 geophysical products meet internationally recognized standards, such as SI traceability, and align with QA4EO principles to guarantee high data quality. Cal/Val aims to validate the accuracy, precision, and stability of L2 products using independent datasets, such as Fiducial Reference Measurements (FRM), and cross\sphinxhyphen{}sensor harmonization. A critical focus is on quantifying and documenting uncertainty propagation from instrument measurements to retrievals while ensuring the retrieval algorithms are robust under diverse environmental conditions.

\sphinxAtStartPar
To meet performance requirements, L2 products must achieve predefined thresholds in accuracy, precision, and stability, with uncertainties traceable and reliable across pixel and aggregate levels. Products must demonstrate consistency over full mission coverage, addressing global, regional, and temporal scales, and maintain defined operational latencies (e.g., near\sphinxhyphen{}real\sphinxhyphen{}time or reprocessed outputs). Validation activities use measurable metrics, including statistical bias, RMS error, and spatial\sphinxhyphen{}temporal coverage rates, to assess product quality systematically. Lastly, the outputs should be delivered in user\sphinxhyphen{}ready formats with comprehensive documentation while maintaining version control and data provenance.


\subsection{ATBD purpose}
\label{\detokenize{01_introduction:atbd-purpose}}
\sphinxAtStartPar
The aim of the Calibration and Validation Algorithm Theoretical Basis Document (ATBD) is to provide a comprehensive and transparent description of the scientific principles, mathematical methods, and processing steps used to generate ROSE\sphinxhyphen{}L L2 ocean products. It ensures that the algorithms are well\sphinxhyphen{}documented, reproducible, and traceable, facilitating a clear understanding of how measurements are transformed into geophysical quantities. The ATBD serves as a foundation for ensuring data quality, guiding calibration/validation efforts, and supporting user confidence in the satellite mission’s outputs.


\subsection{Applicable and reference documents}
\label{\detokenize{01_introduction:applicable-and-reference-documents}}
\sphinxAtStartPar
The source contains no separate applicable\sphinxhyphen{}document list. References cited in the uncertainty\sphinxhyphen{}analysis material are listed in that chapter, and the complete source BibTeX database is included as \sphinxcode{\sphinxupquote{bibliography.bib}}.


\subsection{Acronyms}
\label{\detokenize{01_introduction:acronyms}}\begin{description}
\sphinxlineitem{CANVAS\index{CANVAS@\spxentry{CANVAS}|spxpagem}\phantomsection\label{\detokenize{01_introduction:term-CANVAS}}}
\sphinxAtStartPar
CAlibration aNd VAlidation for Sar.

\sphinxlineitem{CEOS\index{CEOS@\spxentry{CEOS}|spxpagem}\phantomsection\label{\detokenize{01_introduction:term-CEOS}}}
\sphinxAtStartPar
Committee on Earth Observation Satellites.

\sphinxlineitem{DC\index{DC@\spxentry{DC}|spxpagem}\phantomsection\label{\detokenize{01_introduction:term-DC}}}
\sphinxAtStartPar
Doppler Centroid.

\sphinxlineitem{DCA\index{DCA@\spxentry{DCA}|spxpagem}\phantomsection\label{\detokenize{01_introduction:term-DCA}}}
\sphinxAtStartPar
Doppler Centroid Anomaly.

\sphinxlineitem{DORIS\index{DORIS@\spxentry{DORIS}|spxpagem}\phantomsection\label{\detokenize{01_introduction:term-DORIS}}}
\sphinxAtStartPar
Doppler Orbitography and Radiopositioning Integrated by Satellite.

\sphinxlineitem{EO\index{EO@\spxentry{EO}|spxpagem}\phantomsection\label{\detokenize{01_introduction:term-EO}}}
\sphinxAtStartPar
Earth observation.

\sphinxlineitem{ESA\index{ESA@\spxentry{ESA}|spxpagem}\phantomsection\label{\detokenize{01_introduction:term-ESA}}}
\sphinxAtStartPar
European Space Agency.

\sphinxlineitem{EU\index{EU@\spxentry{EU}|spxpagem}\phantomsection\label{\detokenize{01_introduction:term-EU}}}
\sphinxAtStartPar
European Union.

\sphinxlineitem{FDR\index{FDR@\spxentry{FDR}|spxpagem}\phantomsection\label{\detokenize{01_introduction:term-FDR}}}
\sphinxAtStartPar
Fundamental Data Record.

\sphinxlineitem{FIDUCEO\index{FIDUCEO@\spxentry{FIDUCEO}|spxpagem}\phantomsection\label{\detokenize{01_introduction:term-FIDUCEO}}}
\sphinxAtStartPar
Fidelity and Uncertainty in Climate Data Records from Earth Observation.

\sphinxlineitem{FP7\index{FP7@\spxentry{FP7}|spxpagem}\phantomsection\label{\detokenize{01_introduction:term-FP7}}}
\sphinxAtStartPar
Seventh Framework Programme.

\sphinxlineitem{FRM\index{FRM@\spxentry{FRM}|spxpagem}\phantomsection\label{\detokenize{01_introduction:term-FRM}}}
\sphinxAtStartPar
Fiducial Reference Measurement.

\sphinxlineitem{GEO\index{GEO@\spxentry{GEO}|spxpagem}\phantomsection\label{\detokenize{01_introduction:term-GEO}}}
\sphinxAtStartPar
Group on Earth Observations.

\sphinxlineitem{GEOSS\index{GEOSS@\spxentry{GEOSS}|spxpagem}\phantomsection\label{\detokenize{01_introduction:term-GEOSS}}}
\sphinxAtStartPar
Global Earth Observation System of Systems.

\sphinxlineitem{GUM\index{GUM@\spxentry{GUM}|spxpagem}\phantomsection\label{\detokenize{01_introduction:term-GUM}}}
\sphinxAtStartPar
Guide to the expression of Uncertainty in Measurement.

\sphinxlineitem{JCGM\index{JCGM@\spxentry{JCGM}|spxpagem}\phantomsection\label{\detokenize{01_introduction:term-JCGM}}}
\sphinxAtStartPar
Joint Committee for Guides in Metrology.

\sphinxlineitem{QA4EO\index{QA4EO@\spxentry{QA4EO}|spxpagem}\phantomsection\label{\detokenize{01_introduction:term-QA4EO}}}
\sphinxAtStartPar
Quality Assurance Framework for Earth Observation.

\sphinxlineitem{SAR\index{SAR@\spxentry{SAR}|spxpagem}\phantomsection\label{\detokenize{01_introduction:term-SAR}}}
\sphinxAtStartPar
Synthetic Aperture Radar.

\sphinxlineitem{SI\index{SI@\spxentry{SI}|spxpagem}\phantomsection\label{\detokenize{01_introduction:term-SI}}}
\sphinxAtStartPar
International System of Units.

\sphinxlineitem{TDP\index{TDP@\spxentry{TDP}|spxpagem}\phantomsection\label{\detokenize{01_introduction:term-TDP}}}
\sphinxAtStartPar
Thematic Data Product.

\sphinxlineitem{UTD\index{UTD@\spxentry{UTD}|spxpagem}\phantomsection\label{\detokenize{01_introduction:term-UTD}}}
\sphinxAtStartPar
Uncertainty Tree Diagram.

\sphinxlineitem{VIM\index{VIM@\spxentry{VIM}|spxpagem}\phantomsection\label{\detokenize{01_introduction:term-VIM}}}
\sphinxAtStartPar
International Vocabulary of Metrology.

\end{description}


\subsection{Definitions and conventions}
\label{\detokenize{01_introduction:definitions-and-conventions}}\begin{description}
\sphinxlineitem{Accuracy\index{Accuracy@\spxentry{Accuracy}|spxpagem}\phantomsection\label{\detokenize{01_introduction:term-Accuracy}}}
\sphinxAtStartPar
The quality of being near to the true value. The source marks this
definition as requiring a reference.

\sphinxlineitem{Comparison\index{Comparison@\spxentry{Comparison}|spxpagem}\phantomsection\label{\detokenize{01_introduction:term-Comparison}}}
\sphinxAtStartPar
Metrologists validate uncertainty analysis and confirm traceability
through comparisons. The Mutual Recognition Arrangement requires
regular, formal international comparisons to establish the degree of
equivalence between SI realisations and disseminations in different
countries.

\sphinxlineitem{Dataset\index{Dataset@\spxentry{Dataset}|spxpagem}\phantomsection\label{\detokenize{01_introduction:term-Dataset}}}
\sphinxAtStartPar
A combination of data records, such as numerical values needed to analyse
and interpret a natural process, and the related information content.

\sphinxlineitem{Fiducial Reference Measurement\index{Fiducial Reference Measurement@\spxentry{Fiducial Reference Measurement}|spxpagem}\phantomsection\label{\detokenize{01_introduction:term-Fiducial-Reference-Measurement}}}
\sphinxAtStartPar
A suite of independent, fully characterised, and traceable sub\sphinxhyphen{}orbital
measurements, relevant to space\sphinxhyphen{}based observations, following the
guidelines of the GEO/CEOS Quality Assurance Framework for Earth
Observation (QA4EO).

\sphinxlineitem{Fundamental Data Record\index{Fundamental Data Record@\spxentry{Fundamental Data Record}|spxpagem}\phantomsection\label{\detokenize{01_introduction:term-Fundamental-Data-Record}}}
\sphinxAtStartPar
A record, of sufficient duration for its application, of
uncertainty\sphinxhyphen{}quantified sensor observations calibrated to physical units
and located in time and space, together with ancillary and lower\sphinxhyphen{}level
instrument data used to calibrate and locate the observations and to
estimate uncertainty ({\hyperref[\detokenize{04_uncertainty_analysis:ref-qa4eo1}]{\sphinxcrossref{\DUrole{std}{\DUrole{std-ref}{Woolliams et al. 2025a}}}}}).

\sphinxlineitem{Interoperability\index{Interoperability@\spxentry{Interoperability}|spxpagem}\phantomsection\label{\detokenize{01_introduction:term-Interoperability}}}
\sphinxAtStartPar
The ability of data or tools from non\sphinxhyphen{}cooperating resources to integrate
or work together with minimal effort (\sphinxhref{https://doi.org/10.1038/sdata.2016.18}{Wilkinson et al. (2016)}%
\begin{footnote}[2]\sphinxAtStartFootnote
\sphinxnolinkurl{https://doi.org/10.1038/sdata.2016.18}
%
\end{footnote}).

\sphinxlineitem{Interoperable\index{Interoperable@\spxentry{Interoperable}|spxpagem}\phantomsection\label{\detokenize{01_introduction:term-Interoperable}}}
\sphinxAtStartPar
Describes data or tools from non\sphinxhyphen{}cooperating resources that can integrate
or work together with minimal effort (\sphinxhref{https://doi.org/10.1038/sdata.2016.18}{Wilkinson et al. (2016)}%
\begin{footnote}[3]\sphinxAtStartFootnote
\sphinxnolinkurl{https://doi.org/10.1038/sdata.2016.18}
%
\end{footnote}).

\sphinxlineitem{Measurand\index{Measurand@\spxentry{Measurand}|spxpagem}\phantomsection\label{\detokenize{01_introduction:term-Measurand}}}
\sphinxAtStartPar
The quantity intended to be measured, and the single output quantity from
a measurement model.

\sphinxlineitem{Metrology\index{Metrology@\spxentry{Metrology}|spxpagem}\phantomsection\label{\detokenize{01_introduction:term-Metrology}}}
\sphinxAtStartPar
The science of measurement and its application
({\hyperref[\detokenize{04_uncertainty_analysis:ref-jcgmvim}]{\sphinxcrossref{\DUrole{std}{\DUrole{std-ref}{BIPM et al., n.d.\sphinxhyphen{}b}}}}}).

\sphinxlineitem{Thematic Data Product\index{Thematic Data Product@\spxentry{Thematic Data Product}|spxpagem}\phantomsection\label{\detokenize{01_introduction:term-Thematic-Data-Product}}}
\sphinxAtStartPar
A record, of sufficient duration for its application, of
uncertainty\sphinxhyphen{}quantified retrieved values of a geophysical variable, along
with ancillary data used in retrieval and uncertainty estimation.

\sphinxlineitem{Traceability\index{Traceability@\spxentry{Traceability}|spxpagem}\phantomsection\label{\detokenize{01_introduction:term-Traceability}}}
\sphinxAtStartPar
A property of a measurement relating its value to a stated metrological
reference through an unbroken chain of calibrations or comparisons, with
uncertainties evaluated and methodologies documented at each step.

\sphinxlineitem{Uncertainty analysis\index{Uncertainty analysis@\spxentry{Uncertainty analysis}|spxpagem}\phantomsection\label{\detokenize{01_introduction:term-Uncertainty-analysis}}}
\sphinxAtStartPar
The review of uncertainty sources and their propagation through the
traceability chain, based on the suite of documents prepared by the JCGM
and known as the GUM.

\end{description}

\sphinxstepscope


\section{Processing model and data flow}
\label{\detokenize{02_processing_model:processing-model-and-data-flow}}\label{\detokenize{02_processing_model::doc}}

\subsection{Overview and top\sphinxhyphen{}down decomposition}
\label{\detokenize{02_processing_model:overview-and-top-down-decomposition}}
\sphinxAtStartPar
The source describes the calibration and validation activities across the
ROSE\sphinxhyphen{}L mission phases. It does not provide a detailed processing\sphinxhyphen{}block
decomposition or an algorithm\sphinxhyphen{}level data\sphinxhyphen{}flow diagram.


\subsection{Pre\sphinxhyphen{}launch phase (D)}
\label{\detokenize{02_processing_model:pre-launch-phase-d}}
\sphinxAtStartPar
During the pre\sphinxhyphen{}launch phase (phase D), activities focus on preparing and
validating the L2 processing chain, ensuring readiness for on\sphinxhyphen{}orbit operations.
Key tasks include finalizing L2 requirements, algorithms, and performance
budgets. Simulators and proxy datasets are leveraged to validate retrieval
algorithms and uncertainty models. Pre\sphinxhyphen{}launch instrument characterization is
integrated into the processing chain to ensure retrievals are robust to on\sphinxhyphen{}orbit
conditions. A comprehensive Cal/Val Plan is developed, identifying validation
sites, protocols, data matchups, and metrics.


\subsection{Commissioning phase (IOC or E1)}
\label{\detokenize{02_processing_model:commissioning-phase-ioc-or-e1}}
\sphinxAtStartPar
In the commissioning phase (IOC or phase E1), the primary goal is to validate
L2 products using on\sphinxhyphen{}orbit data and release initial beta/provisional products
with documented performance. Early activities include confirming L1
performance, performing vicarious calibrations, and matching L2 products with
ground\sphinxhyphen{}based and airborne reference datasets (FRM sites, buoys, ship\sphinxhyphen{}based
campaigns, etc.). Algorithms are fine\sphinxhyphen{}tuned using real\sphinxhyphen{}time data, and initial
uncertainty models are validated. Early reprocessing cycles are conducted to
improve product stability, and performance against key metrics is assessed.
Products are released with quality summaries, and ongoing monitoring dashboards
are set up.


\subsection{Exploitation phase (E2)}
\label{\detokenize{02_processing_model:exploitation-phase-e2}}
\sphinxAtStartPar
The exploitation phase (phase E2) ensures L2 products maintain accuracy,
precision, and stability over the mission lifecycle. Routine validation is
conducted using reference networks, continuous KPI monitoring, and periodic
reprocessing cycles for algorithm and calibration improvements. Traceability
to standards (QA4EO, FRM) and uncertainty validation is maintained, while
long\sphinxhyphen{}term drift corrections support stable time series. Community engagement,
user feedback, and updates to documentation (validation reports, product
guides) are prioritised. The goal is to achieve validated product maturity,
ensuring high\sphinxhyphen{}quality data for science and operational applications.


\subsection{End\sphinxhyphen{}to\sphinxhyphen{}end data flow}
\label{\detokenize{02_processing_model:end-to-end-data-flow}}
\sphinxAtStartPar
No end\sphinxhyphen{}to\sphinxhyphen{}end processing\sphinxhyphen{}flow diagram, processing\sphinxhyphen{}step mapping, or algorithm
identifiers are provided in the source LaTeX document.

\sphinxstepscope


\section{Algorithms description}
\label{\detokenize{03_algorithms:algorithms-description}}\label{\detokenize{03_algorithms::doc}}
\sphinxAtStartPar
The source groups its method outlines into calibration and validation. The
method\sphinxhyphen{}specific headings are retained in the following pages. Their details
are not present in the source.

\sphinxstepscope


\subsection{Selected Calibration methods}
\label{\detokenize{alg-01-calibration:selected-calibration-methods}}\label{\detokenize{alg-01-calibration::doc}}
\sphinxAtStartPar
The source provides the following section headings but no method descriptions,
input specifications, processing details, performance metrics, or expected
outputs.


\subsubsection{ROSE\sphinxhyphen{}L L1 SLC products}
\label{\detokenize{alg-01-calibration:rose-l-l1-slc-products}}

\paragraph{Short description}
\label{\detokenize{alg-01-calibration:short-description}}

\paragraph{Description of inputs}
\label{\detokenize{alg-01-calibration:description-of-inputs}}

\paragraph{Detailed methodology}
\label{\detokenize{alg-01-calibration:detailed-methodology}}

\paragraph{Performance metrics}
\label{\detokenize{alg-01-calibration:performance-metrics}}

\paragraph{Expected output}
\label{\detokenize{alg-01-calibration:expected-output}}

\subsubsection{NOC}
\label{\detokenize{alg-01-calibration:noc}}

\paragraph{Short description}
\label{\detokenize{alg-01-calibration:id1}}

\paragraph{Description of inputs}
\label{\detokenize{alg-01-calibration:id2}}

\paragraph{Detailed methodology}
\label{\detokenize{alg-01-calibration:id3}}

\paragraph{Performance metrics}
\label{\detokenize{alg-01-calibration:id4}}

\paragraph{Expected output}
\label{\detokenize{alg-01-calibration:id5}}

\subsubsection{GMF tuning}
\label{\detokenize{alg-01-calibration:gmf-tuning}}

\paragraph{Short description}
\label{\detokenize{alg-01-calibration:id6}}

\paragraph{Description of inputs}
\label{\detokenize{alg-01-calibration:id7}}

\paragraph{Detailed methodology}
\label{\detokenize{alg-01-calibration:id8}}

\paragraph{Performance metrics}
\label{\detokenize{alg-01-calibration:id9}}

\paragraph{Expected output}
\label{\detokenize{alg-01-calibration:id10}}

\subsubsection{MTF tuning}
\label{\detokenize{alg-01-calibration:mtf-tuning}}

\paragraph{Short description}
\label{\detokenize{alg-01-calibration:id11}}

\paragraph{Description of inputs}
\label{\detokenize{alg-01-calibration:id12}}

\paragraph{Detailed methodology}
\label{\detokenize{alg-01-calibration:id13}}

\paragraph{Performance metrics}
\label{\detokenize{alg-01-calibration:id14}}

\paragraph{Expected output}
\label{\detokenize{alg-01-calibration:id15}}

\subsubsection{Summary}
\label{\detokenize{alg-01-calibration:summary}}
\sphinxAtStartPar
The source contains this heading without summary text.

\sphinxstepscope


\subsection{Selected Validation methods}
\label{\detokenize{alg-02-validation:selected-validation-methods}}\label{\detokenize{alg-02-validation::doc}}
\sphinxAtStartPar
The source provides the following section headings but no method descriptions,
input specifications, processing details, performance metrics, or expected
outputs.


\subsubsection{One\sphinxhyphen{}to\sphinxhyphen{}one validation}
\label{\detokenize{alg-02-validation:one-to-one-validation}}

\paragraph{Short description}
\label{\detokenize{alg-02-validation:short-description}}

\paragraph{Description of inputs}
\label{\detokenize{alg-02-validation:description-of-inputs}}

\paragraph{Detailed methodology}
\label{\detokenize{alg-02-validation:detailed-methodology}}

\paragraph{Performance metrics}
\label{\detokenize{alg-02-validation:performance-metrics}}

\paragraph{Expected output}
\label{\detokenize{alg-02-validation:expected-output}}

\subsubsection{Triple collocation}
\label{\detokenize{alg-02-validation:triple-collocation}}

\paragraph{Short description}
\label{\detokenize{alg-02-validation:id1}}

\paragraph{Description of inputs}
\label{\detokenize{alg-02-validation:id2}}

\paragraph{Detailed methodology}
\label{\detokenize{alg-02-validation:id3}}

\paragraph{Performance metrics}
\label{\detokenize{alg-02-validation:id4}}

\paragraph{Expected output}
\label{\detokenize{alg-02-validation:id5}}

\subsubsection{Spatial variances}
\label{\detokenize{alg-02-validation:spatial-variances}}

\paragraph{Short description}
\label{\detokenize{alg-02-validation:id6}}

\paragraph{Description of inputs}
\label{\detokenize{alg-02-validation:id7}}

\paragraph{Detailed methodology}
\label{\detokenize{alg-02-validation:id8}}

\paragraph{Performance metrics}
\label{\detokenize{alg-02-validation:id9}}

\paragraph{Expected output}
\label{\detokenize{alg-02-validation:id10}}

\subsubsection{Summary}
\label{\detokenize{alg-02-validation:summary}}
\sphinxAtStartPar
The source contains this heading without summary text.

\sphinxstepscope


\section{Uncertainty analysis}
\label{\detokenize{04_uncertainty_analysis:uncertainty-analysis}}\label{\detokenize{04_uncertainty_analysis::doc}}

\subsection{Introduction}
\label{\detokenize{04_uncertainty_analysis:introduction}}
\sphinxAtStartPar
The “Evaluation of measurement data \textendash{} GUM” {\hyperref[\detokenize{04_uncertainty_analysis:ref-jcgmgum}]{\sphinxcrossref{\DUrole{std}{\DUrole{std-ref}{(BIPM et al., n.d.\sphinxhyphen{}a)}}}}}, or the “GUM bible”, is a foundational international document that provides a conceptual and mathematical framework for identifying sources of uncertainty, quantifying them in a consistent way (as standard deviations), combining them into a single uncertainty figure, and reporting the result and its uncertainty clearly and comparably across disciplines and countries. Focusing on evaluation and expression of uncertainty in measurement it is an important component of the science of measurement and application, known as metrology (e.g., {\hyperref[\detokenize{04_uncertainty_analysis:ref-jcgmvim}]{\sphinxcrossref{\DUrole{std}{\DUrole{std-ref}{BIPM et al., n.d.\sphinxhyphen{}b}}}}}).

\sphinxAtStartPar
The GEO is an international partnership of governments and organizations acting as a policy and coordination body to promote global sharing and use of EO data, for example to support work on climate, disaster management, agriculture, water, and ecosystems. The GEOSS is the global, distributed infrastructure that GEO is building and coordinating. Rather than being a single system, it connects many independent EO systems \textendash{} such as satellites, ground networks, and ocean buoys \textendash{} from different countries and agencies, and aims to make them interoperable, accessible through common standards and interfaces, and accompanied by information about data quality and uncertainty. As such, GEO is the partnership, and GEOSS is the “system of systems” that GEO oversees.

\sphinxAtStartPar
To support its mission, GEOSS needs a metrological framework. As such, the 2010 endorsement of the QA4EO by the CEOS (part of GEOSS), has defined guidelines for EO applications based on principles adapted from the metrology community {\hyperref[\detokenize{04_uncertainty_analysis:ref-qa4eo1}]{\sphinxcrossref{\DUrole{std}{\DUrole{std-ref}{(Woolliams et al. 2025a)}}}}}. QA4EO aims to ensure that EO data are accessible, of known and documented quality, traceable to reference standards, in particular the SI, and suitable for user application (“fit for purpose”).

\sphinxAtStartPar
This section applies the GUM principles to SAR data and reviews the existing frameworks supporting metrology: the VIM {\hyperref[\detokenize{04_uncertainty_analysis:ref-jcgmvim}]{\sphinxcrossref{\DUrole{std}{\DUrole{std-ref}{(BIPM et al., n.d.\sphinxhyphen{}b)}}}}}, the GUM {\hyperref[\detokenize{04_uncertainty_analysis:ref-jcgmgum}]{\sphinxcrossref{\DUrole{std}{\DUrole{std-ref}{(BIPM et al., n.d.\sphinxhyphen{}a)}}}}}, and QA4EO. The LaTeX source also refers to section labels for level\sphinxhyphen{}1, level\sphinxhyphen{}2, error propagation, and conclusions that are not defined there. Its product\sphinxhyphen{}specific uncertainty\sphinxhyphen{}estimation headings are retained below.


\subsection{The GUM Principles: A review}
\label{\detokenize{04_uncertainty_analysis:the-gum-principles-a-review}}

\subsubsection{The International Vocabulary of Metrology (VIM)}
\label{\detokenize{04_uncertainty_analysis:the-international-vocabulary-of-metrology-vim}}
\sphinxAtStartPar
The International Vocabulary of Metrology (VIM; {\hyperref[\detokenize{04_uncertainty_analysis:ref-jcgmvim}]{\sphinxcrossref{\DUrole{std}{\DUrole{std-ref}{BIPM et al., n.d.\sphinxhyphen{}b}}}}}) defines and organizes key concepts in metrology to ensure global consistency in measurement science. It provides standardized terminology covering quantities, units, and measuring instruments, with a focus on measurement uncertainty rather than error.

\sphinxAtStartPar
The document explains critical properties of measuring devices, such as precision, stability, and sensitivity, and highlights the importance of traceability and standards for reliability.

\sphinxAtStartPar
Its applications span various fields, including physics, chemistry, biology, medicine, and emerging areas like forensic science and food technology. Harmonization of terms is emphasized to ensure coherence in the various scientific, engineering, regulatory, and quality assurance fields.


\subsubsection{Guide to the Expression of Uncertainty in Measurement (GUM)}
\label{\detokenize{04_uncertainty_analysis:guide-to-the-expression-of-uncertainty-in-measurement-gum}}
\sphinxAtStartPar
The GUM establishes general rules for evaluating and expressing uncertainty in measurement that are intended to be applicable to a broad spectrum of measurements {\hyperref[\detokenize{04_uncertainty_analysis:ref-jcgmgum}]{\sphinxcrossref{\DUrole{std}{\DUrole{std-ref}{(BIPM et al., n.d.\sphinxhyphen{}a)}}}}}. The guide was developed to solve the challenge of different laboratories and disciplines using different methods and terminology for uncertainty, which made measurement results difficult to compare.

\sphinxAtStartPar
The central idea of the GUM is to treat all components of uncertainty in the same mathematical way, regardless of whether they arise from random variation or systematic effects. Every component is represented by a probability distribution and a standard uncertainty. To obtain standard uncertainties, the GUM distinguishes between two kinds of evaluation. “Type A” evaluations use statistical analysis of repeated measurements, whereas “Type B” evaluations use all other available information, such as instrument specifications, calibration certificates, or expert judgment. The results of both evaluation types can be converted into standard uncertainties so they can be used on the same terms.

\sphinxAtStartPar
Measurements are described by a mathematical model that links the quantity being measured (the measurand) to various input, or influence, quantities. Once standard uncertainties of the influence quantities are estimated, they are combined \textendash{} taking into account any correlations \textendash{} using standard rules for propagating variances. This yields the combined standard uncertainty of the measurand, taken as the estimated standard deviation associated with the result. An expanded uncertainty is obtained by multiplying the standard uncertainty by a coverage factor \(k\). The purpose of the expanded uncertainty is to provide an interval about the result of a measurement that may be expected to encompass a large fraction of the distribution of values that could reasonably be attributed to the measurand. The value of \(k\) and how it was chosen must always be stated.

\sphinxAtStartPar
A key emphasis of the GUM is transparency. An uncertainty statement should clearly describe which inputs were considered, how each uncertainty component was obtained, what assumptions were made, and how the components were combined. The document supplements this framework with definitions, statistical background, guidance on practical evaluation, and worked examples. In doing so, it provides a common language and method that allow measurement results and their uncertainties to be understood and compared across borders and domains.


\subsubsection{Quality Assurance Framework for Earth Observation (QA4EO)}
\label{\detokenize{04_uncertainty_analysis:quality-assurance-framework-for-earth-observation-qa4eo}}
\sphinxAtStartPar
GEOSS depends on interoperable and quality\sphinxhyphen{}assured data to derive comprehensive information about climate and environment, as well as, e.g., food, water and energy security. To achieve interoperability across systems and datasets, common and well\sphinxhyphen{}defined data and metadata formats are needed, but also metrological principles such as traceability to SI, documented uncertainty analysis and comparison between different measurement references must be followed, and are key to the QA4EO reports {\hyperref[\detokenize{04_uncertainty_analysis:ref-qa4eo1}]{\sphinxcrossref{\DUrole{std}{\DUrole{std-ref}{(Woolliams et al. 2025a)}}}}}.

\sphinxAtStartPar
Beside {\hyperref[\detokenize{04_uncertainty_analysis:ref-jcgmvim}]{\sphinxcrossref{\DUrole{std}{\DUrole{std-ref}{BIPM et al., n.d.\sphinxhyphen{}b}}}}} and {\hyperref[\detokenize{04_uncertainty_analysis:ref-jcgmgum}]{\sphinxcrossref{\DUrole{std}{\DUrole{std-ref}{BIPM et al., n.d.\sphinxhyphen{}a}}}}}, QA4EO also builds on EO\sphinxhyphen{}metrology collaboration in the EU FP7 and Horizon 2020, ESA and Copernicus projects and, in particular, the FIDUCEO project {\hyperref[\detokenize{04_uncertainty_analysis:ref-mittaz19}]{\sphinxcrossref{\DUrole{std}{\DUrole{std-ref}{(Mittaz et al. 2019)}}}}}.

\sphinxAtStartPar
As summarized in {\hyperref[\detokenize{04_uncertainty_analysis:ref-qa4eo2}]{\sphinxcrossref{\DUrole{std}{\DUrole{std-ref}{(Woolliams et al. 2025b)}}}}}, three basic metrology principles generally apply to EO data:
\begin{enumerate}
\sphinxsetlistlabels{\arabic}{enumi}{enumii}{}{.}%
\item {} 
\sphinxAtStartPar
\sphinxstylestrong{Traceability:} Metrological traceability is a property of a measurement that relates the measured value to a stated metrological reference through an unbroken chain of calibrations or comparisons. It requires, for each step in the traceability chain, that uncertainties are evaluated and that methodologies are documented.

\item {} 
\sphinxAtStartPar
\sphinxstylestrong{Uncertainty Analysis:} Uncertainty analysis is the review of all sources of uncertainty and the propagation of that uncertainty through the traceability chain. It is based on the suite of documents prepared by the JCGM and together known as the GUM.

\item {} 
\sphinxAtStartPar
\sphinxstylestrong{Comparison:} Metrologists validate uncertainty analysis and confirm traceability through comparisons. The Mutual Recognition Arrangement (a formal arrangement to confirm the degree\sphinxhyphen{}of\sphinxhyphen{}equivalence between SI realisations and disseminations in different countries) requires regular, formal international comparisons that are conducted under strict rules.

\end{enumerate}

\sphinxAtStartPar
The uncertainty analysis, as outlined in the GUM, should follow these principles.

\sphinxAtStartPar
ESA has initially applied the terms FRM, FDR and TDP to describe metrologically rigorous observations of specific relevance to space\sphinxhyphen{}based observations {\hyperref[\detokenize{04_uncertainty_analysis:ref-qa4eo2}]{\sphinxcrossref{\DUrole{std}{\DUrole{std-ref}{(Woolliams et al. 2025b)}}}}}. The terms are increasingly being used by the broader EO community, although they have not yet been formally endorsed. The terms are defined in {\hyperref[\detokenize{04_uncertainty_analysis:ref-qa4eo2}]{\sphinxcrossref{\DUrole{std}{\DUrole{std-ref}{Woolliams et al. 2025b}}}}}, as follows:
\begin{itemize}
\item {} 
\sphinxAtStartPar
\sphinxstylestrong{Fundamental Data Record (FDR):} An FDR is a record, of sufficient duration for its application, of uncertainty\sphinxhyphen{}quantified sensor observations calibrated to physical units and located in time and space, together with all ancillary and lower\sphinxhyphen{}level instrument data used to calibrate and locate the observations and to estimate uncertainty.

\item {} 
\sphinxAtStartPar
\sphinxstylestrong{Thematic Data Product (TDP):} A TDP is a record, of sufficient duration for its application, of uncertainty\sphinxhyphen{}quantified retrieved values of a geophysical variable, along with all ancillary data used in retrieval and uncertainty estimation.

\item {} 
\sphinxAtStartPar
\sphinxstylestrong{Fiducial Reference Measurement (FRM):} An FRM is a suite of independent, fully characterised, and traceable (to a community agreed reference, ideally SI) measurements of a satellite relevant measurand, tailored specifically to address the calibration/validation needs of a class of satellite borne sensors and that follow the guidelines outlined by the GEO/CEOS QA4EO.

\end{itemize}

\sphinxAtStartPar
As highlighted by {\hyperref[\detokenize{04_uncertainty_analysis:ref-qa4eo2}]{\sphinxcrossref{\DUrole{std}{\DUrole{std-ref}{Woolliams et al. 2025b}}}}}, “TDPs provide higher level products that have been processed from FDRs, through algorithms which also often combine information from other FDRs (e.g., from other satellite sensors) or from external information (such as reanalysis models and/or certain non\sphinxhyphen{}satellite data), along with such additional information”. Guidelines for assessing observational systems with respect to FRMs have been established by CEOS.

\sphinxAtStartPar
In order to establish a systematic process for generating uncertainty budgets, and ensuring that EO data are interoperable, temporally stable, and fit for scientific, operational and decision\sphinxhyphen{}making purposes, the QA4EO and EO projects acknowledged in {\hyperref[\detokenize{04_uncertainty_analysis:ref-qa4eo3}]{\sphinxcrossref{\DUrole{std}{\DUrole{std-ref}{Woolliams et al. 2025c}}}}} have adapted the stages, suggested in GUM, required to support the development of a measurement model and from that an uncertainty analysis. These EO specific stages of uncertainty analysis are given by {\hyperref[\detokenize{04_uncertainty_analysis:ref-qa4eo3}]{\sphinxcrossref{\DUrole{std}{\DUrole{std-ref}{Woolliams et al. 2025c}}}}} as follows:
\begin{itemize}
\item {} 
\sphinxAtStartPar
Step 1: Define the measurand and measurement model.

\item {} 
\sphinxAtStartPar
Step 2: Establish the traceability with a diagram.

\item {} 
\sphinxAtStartPar
Step 3: Evaluate each source of uncertainty and fill out an effects table.

\item {} 
\sphinxAtStartPar
Step 4: Calculate the data product and uncertainties.

\item {} 
\sphinxAtStartPar
Step 5: Record information about the uncertainty analysis for long term data preservation purposes and summarise for today’s users.

\end{itemize}

\sphinxAtStartPar
These steps are discussed in detail in {\hyperref[\detokenize{04_uncertainty_analysis:ref-qa4eo3}]{\sphinxcrossref{\DUrole{std}{\DUrole{std-ref}{Woolliams et al. 2025c}}}}}. The product\sphinxhyphen{}specific headings below retain the source’s intended level\sphinxhyphen{}1 and level\sphinxhyphen{}2 uncertainty\sphinxhyphen{}estimation scope.


\subsection{Product\sphinxhyphen{}specific uncertainty estimation}
\label{\detokenize{04_uncertainty_analysis:product-specific-uncertainty-estimation}}
\sphinxAtStartPar
The source provides headings for the following product analyses, but no detailed
measurement models, traceability matrices, effect evaluations, uncertainty
calculations, or output documentation beneath them.


\subsubsection{Azimuth wavelength cutoff}
\label{\detokenize{04_uncertainty_analysis:azimuth-wavelength-cutoff}}\begin{itemize}
\item {} 
\sphinxAtStartPar
Description of measurand and measurement model

\item {} 
\sphinxAtStartPar
Traceability matrix

\item {} 
\sphinxAtStartPar
Evaluation of effects

\item {} 
\sphinxAtStartPar
Calculation of associated uncertainties

\item {} 
\sphinxAtStartPar
Documentation and outputs

\end{itemize}


\subsubsection{MACS}
\label{\detokenize{04_uncertainty_analysis:macs}}\begin{itemize}
\item {} 
\sphinxAtStartPar
Description of measurand and measurement model

\item {} 
\sphinxAtStartPar
Traceability matrix

\item {} 
\sphinxAtStartPar
Evaluation of effects

\item {} 
\sphinxAtStartPar
Calculation of associated uncertainties

\item {} 
\sphinxAtStartPar
Documentation and outputs

\end{itemize}


\subsubsection{iMACS}
\label{\detokenize{04_uncertainty_analysis:imacs}}\begin{itemize}
\item {} 
\sphinxAtStartPar
Description of measurand and measurement model

\item {} 
\sphinxAtStartPar
Traceability matrix

\item {} 
\sphinxAtStartPar
Evaluation of effects

\item {} 
\sphinxAtStartPar
Calculation of associated uncertainties

\item {} 
\sphinxAtStartPar
Documentation and outputs

\end{itemize}


\subsubsection{DCA}
\label{\detokenize{04_uncertainty_analysis:dca}}\begin{itemize}
\item {} 
\sphinxAtStartPar
Description of measurand and measurement model

\item {} 
\sphinxAtStartPar
Traceability matrix

\item {} 
\sphinxAtStartPar
Evaluation of effects

\item {} 
\sphinxAtStartPar
Calculation of associated uncertainties

\item {} 
\sphinxAtStartPar
Documentation and outputs

\end{itemize}


\subsubsection{Texture\sphinxhyphen{}based azimuth cutoff}
\label{\detokenize{04_uncertainty_analysis:texture-based-azimuth-cutoff}}\begin{itemize}
\item {} 
\sphinxAtStartPar
Description of measurand and measurement model

\item {} 
\sphinxAtStartPar
Traceability matrix

\item {} 
\sphinxAtStartPar
Evaluation of effects

\item {} 
\sphinxAtStartPar
Calculation of associated uncertainties

\item {} 
\sphinxAtStartPar
Documentation and outputs

\end{itemize}


\subsubsection{RVL}
\label{\detokenize{04_uncertainty_analysis:rvl}}\begin{itemize}
\item {} 
\sphinxAtStartPar
Description of measurand and measurement model

\item {} 
\sphinxAtStartPar
Traceability matrix

\item {} 
\sphinxAtStartPar
Evaluation of effects

\item {} 
\sphinxAtStartPar
Calculation of associated uncertainties

\item {} 
\sphinxAtStartPar
Documentation and outputs

\end{itemize}


\subsubsection{Inverse Wave Age}
\label{\detokenize{04_uncertainty_analysis:inverse-wave-age}}\begin{itemize}
\item {} 
\sphinxAtStartPar
Description of measurand and measurement model

\item {} 
\sphinxAtStartPar
Traceability matrix

\item {} 
\sphinxAtStartPar
Evaluation of effects

\item {} 
\sphinxAtStartPar
Calculation of associated uncertainties

\item {} 
\sphinxAtStartPar
Documentation and outputs

\end{itemize}


\subsubsection{Wind products}
\label{\detokenize{04_uncertainty_analysis:wind-products}}

\paragraph{Wind speed, wind stress, and hub height wind speed}
\label{\detokenize{04_uncertainty_analysis:wind-speed-wind-stress-and-hub-height-wind-speed}}\begin{itemize}
\item {} 
\sphinxAtStartPar
Description of measurand and measurement model

\item {} 
\sphinxAtStartPar
Traceability matrix

\item {} 
\sphinxAtStartPar
Evaluation of effects

\item {} 
\sphinxAtStartPar
Calculation of associated uncertainties

\item {} 
\sphinxAtStartPar
Documentation and outputs

\end{itemize}


\subsubsection{Sea\sphinxhyphen{}state products}
\label{\detokenize{04_uncertainty_analysis:sea-state-products}}
\sphinxAtStartPar
Total significant wave height; mean period wind\sphinxhyphen{}sea; first and second moment
wave period; wave height for the swell\sphinxhyphen{}dominant and swell\sphinxhyphen{}secondary systems;
significant wave height wind\sphinxhyphen{}sea; wave direction; and wave directional energy
spectrum.
\begin{itemize}
\item {} 
\sphinxAtStartPar
Description of measurand and measurement model

\item {} 
\sphinxAtStartPar
Traceability matrix

\item {} 
\sphinxAtStartPar
Evaluation of effects

\item {} 
\sphinxAtStartPar
Calculation of associated uncertainties

\item {} 
\sphinxAtStartPar
Documentation and outputs

\end{itemize}


\subsubsection{Line\sphinxhyphen{}of\sphinxhyphen{}sight current products}
\label{\detokenize{04_uncertainty_analysis:line-of-sight-current-products}}

\paragraph{Line\sphinxhyphen{}of\sphinxhyphen{}sight wave Doppler and line\sphinxhyphen{}of\sphinxhyphen{}sight surface current component}
\label{\detokenize{04_uncertainty_analysis:line-of-sight-wave-doppler-and-line-of-sight-surface-current-component}}\begin{itemize}
\item {} 
\sphinxAtStartPar
Description of measurand and measurement model

\item {} 
\sphinxAtStartPar
Traceability matrix

\item {} 
\sphinxAtStartPar
Evaluation of effects

\item {} 
\sphinxAtStartPar
Calculation of associated uncertainties

\item {} 
\sphinxAtStartPar
Documentation and outputs

\end{itemize}


\subsubsection{Summary}
\label{\detokenize{04_uncertainty_analysis:summary}}
\sphinxAtStartPar
The source contains this heading without summary text.


\subsection{References}
\label{\detokenize{04_uncertainty_analysis:references}}\phantomsection\label{\detokenize{04_uncertainty_analysis:ref-jcgmgum}}
\sphinxAtStartPar
BIPM, IEC, IFCC, et al. n.d.\sphinxhyphen{}a. \sphinxstyleemphasis{Evaluation of Measurement Data — Guide to the
Expression of Uncertainty in Measurement}. Joint Committee for Guides in
Metrology, JCGM 100:2008. \sphinxhref{https://doi.org/10.59161/JCGM100-2008E}{doi:10.59161/JCGM100\sphinxhyphen{}2008E}%
\begin{footnote}[4]\sphinxAtStartFootnote
\sphinxnolinkurl{https://doi.org/10.59161/JCGM100-2008E}
%
\end{footnote}.

\phantomsection\label{\detokenize{04_uncertainty_analysis:ref-jcgmvim}}
\sphinxAtStartPar
BIPM, IEC, IFCC, et al. n.d.\sphinxhyphen{}b. \sphinxstyleemphasis{International Vocabulary of Metrology \sphinxhyphen{} Basic
and General Concepts and Associated Terms (VIM)}. Joint Committee for Guides in
Metrology, JCGM 200:2012(E/F).

\phantomsection\label{\detokenize{04_uncertainty_analysis:ref-mittaz19}}
\sphinxAtStartPar
Mittaz, Jonathan, Christopher J. Merchant, and Emma R. Woolliams. 2019.
“Applying Principles of Metrology to Historical Earth Observations from
Satellites.” \sphinxstyleemphasis{Metrologia} 56 (3).
\sphinxhref{https://doi.org/10.1088/1681-7575/ab1705}{doi:10.1088/1681\sphinxhyphen{}7575/ab1705}%
\begin{footnote}[5]\sphinxAtStartFootnote
\sphinxnolinkurl{https://doi.org/10.1088/1681-7575/ab1705}
%
\end{footnote}.

\phantomsection\label{\detokenize{04_uncertainty_analysis:ref-qa4eo1}}
\sphinxAtStartPar
Woolliams, Emma, Sajedeh Behnia, Agnieszka Bialek, et al. 2025a. \sphinxstyleemphasis{General
Guidance on a Metrological Approach to Fundamental Data Records (FDR), Thematic
Data Products (TDP) and Fiducial Reference Measurements (FRM) \sphinxhyphen{} Executive
Summary}. National Physical Laboratory.

\phantomsection\label{\detokenize{04_uncertainty_analysis:ref-qa4eo2}}
\sphinxAtStartPar
Woolliams, Emma, Sajedeh Behnia, Agnieszka Bialek, et al. 2025b. \sphinxstyleemphasis{General
Guidance on a Metrological Approach to Fundamental Data Records (FDR), Thematic
Data Products (TDP) and Fiducial Reference Measurements (FRM) \sphinxhyphen{} Metrology
Theoretical Basis}. National Physical Laboratory.

\phantomsection\label{\detokenize{04_uncertainty_analysis:ref-qa4eo3}}
\sphinxAtStartPar
Woolliams, Emma, Sajedeh Behnia, Agnieszka Bialek, et al. 2025c. \sphinxstyleemphasis{General
Guidance on a Metrological Approach to Fundamental Data Records (FDR), Thematic
Data Products (TDP) and Fiducial Reference Measurements (FRM) \sphinxhyphen{} Uncertainty
Analysis Process}. National Physical Laboratory.



\renewcommand{\indexname}{Index}
\printindex
\end{document}